\chapter{2022年全国甲卷（文）}

\section{单选题}

\begin{problem}
已知集合 \(A=\{-2,-1,0,1,2\}\)，\(B=\{x\mid 0\le x<\dfrac{5}{2}\}\)，则 \(A\cap B=\)（\quad）

\choices
    {\(\{0,1,2\}\)}
    {\(\{-2,-1,0\}\)}
    {\(\{0,1\}\)}
    {\(\{1,2\}\)}
\end{problem}

\begin{problem}
某社区通过公益讲座以普及社区居民的垃圾分类知识. 为了解讲座效果，随机抽取 \(10\) 位社区居民，让他们在讲座前和讲座后各回答一份垃圾分类知识问卷，这 \(10\) 位社区居民在讲座前和讲座后问卷答题的正确率如下图：

\bitmapfigure[max width=6.0cm]{img/2022/national_paper_a_liberal/q2_fig1.png}

则（\quad）

\choices
    {讲座前问卷答题的正确率的中位数小于 \(70\%\)}
    {讲座后问卷答题的正确率的平均数大于 \(85\%\)}
    {讲座前问卷答题的正确率的标准差小于讲座后正确率的标准差}
    {讲座后问卷答题的正确率的极差大于讲座前正确率的极差}
\end{problem}

\begin{problem}
若 \(z=1+\i\)，则 \(\left|\i z+3\overline{z}\right|=\)（\quad）

\choices
    {\(4\sqrt{5}\)}
    {\(4\sqrt{2}\)}
    {\(2\sqrt{5}\)}
    {\(2\sqrt{2}\)}
\end{problem}

\begin{problem}
如图，网格纸上绘制的是一个多面体的三视图，网格小正方形的边长为 \(1\)，则该多面体的体积为（\quad）

\FigureLayoutDeclare{img/2022/national_paper_a_liberal/q4_fig1.png}{6.0cm}{21bp 52bp 0bp 4bp}
\bitmapfigure[max width=6.0cm]{img/2022/national_paper_a_liberal/q4_fig1.png}

\choices
    {8}
    {12}
    {16}
    {20}
\end{problem}

\begin{problem}
将函数 \(f(x)=\sin\left(\omega x+\dfrac{\pi}{3}\right)\)（\(\omega>0\)）的图像向左平移 \(\dfrac{\pi}{2}\) 个单位长度后得到曲线 \(C\)，若 \(C\) 关于 \(y\) 轴对称，则 \(\omega\) 的最小值是（\quad）

\choices
    {\(\dfrac{1}{6}\)}
    {\(\dfrac{1}{4}\)}
    {\(\dfrac{1}{3}\)}
    {\(\dfrac{1}{2}\)}
\end{problem}

\begin{problem}
从分别写有 \(1\)，\(2\)，\(3\)，\(4\)，\(5\)，\(6\) 的 \(6\) 张卡片中无放回随机抽取 \(2\) 张，则抽到的 \(2\) 张卡片上的数字之积是 \(4\) 的倍数的概率为（\quad）

\choices
    {\(\dfrac{1}{5}\)}
    {\(\dfrac{1}{3}\)}
    {\(\dfrac{2}{5}\)}
    {\(\dfrac{2}{3}\)}
\end{problem}

\begin{problem}
函数 \(y=\left(3^x-3^{-x}\right)\cos x\) 在区间 \(\left[-\dfrac{\pi}{2},\dfrac{\pi}{2}\right]\) 的图象大致为（\quad）

\choices
    {\choicebitmap[width=0.15\paperwidth]{img/2022/national_paper_a_liberal/q7_optA.png}}
    {\choicebitmap[width=0.15\paperwidth]{img/2022/national_paper_a_liberal/q7_optB.png}}
    {\choicebitmap[width=0.15\paperwidth]{img/2022/national_paper_a_liberal/q7_optC.png}}
    {\choicebitmap[width=0.15\paperwidth]{img/2022/national_paper_a_liberal/q7_optD.png}}
\end{problem}

\begin{problem}
当 \(x=1\) 时，函数 \(f(x)=a\ln x+\dfrac{b}{x}\) 取得最大值 \(-2\)，则 \(f'(2)=\)（\quad）

\choices
    {\(-1\)}
    {\(-\dfrac{1}{2}\)}
    {\(\dfrac{1}{2}\)}
    {1}
\end{problem}

\begin{problem}
在长方体 \(ABCD-A_1B_1C_1D_1\) 中，已知 \(B_1D\) 与平面 \(ABCD\) 和平面 \(AA_1B_1B\) 所成的角均为 \(30^\circ\)，则（\quad）

\bitmapfigure[max width=6.0cm]{img/2022/national_paper_a_liberal/q9_fig1.png}

\choices
    {\(AB=2AD\)}
    {\(AB\) 与平面 \(AB_1C_1D\) 所成的角为 \(30^\circ\)}
    {\(AC=CB_1\)}
    {\(B_1D\) 与平面 \(BB_1C_1C\) 所成的角为 \(45^\circ\)}
\end{problem}

\begin{problem}
甲，乙两个圆锥的母线长相等，侧面展开图的圆心角之和为 \(2\pi\)，侧面积分别为 \(S_{\text{甲}}\) 和 \(S_{\text{乙}}\)，体积分别为 \(V_{\text{甲}}\) 和 \(V_{\text{乙}}\). 若 \(\dfrac{S_{\text{甲}}}{S_{\text{乙}}}=2\)，则 \(\dfrac{V_{\text{甲}}}{V_{\text{乙}}}=\)（\quad）

\choices
    {\(\sqrt{5}\)}
    {\(2\sqrt{2}\)}
    {\(\sqrt{10}\)}
    {\(\dfrac{5\sqrt{10}}{4}\)}
\end{problem}

\begin{problem}
已知椭圆 \(C:\dfrac{x^2}{a^2}+\dfrac{y^2}{b^2}=1\)（\(a>b>0\)）的离心率为 \(\dfrac{1}{3}\)，\(A_1\)，\(A_2\) 分别为 \(C\) 的左，右顶点，\(B\) 为 \(C\) 的上顶点. 若 \(\overrightarrow{BA_1}\cdot\overrightarrow{BA_2}=-1\)，则 \(C\) 的方程为（\quad）

\choices
    {\(\dfrac{x^2}{18}+\dfrac{y^2}{16}=1\)}
    {\(\dfrac{x^2}{9}+\dfrac{y^2}{8}=1\)}
    {\(\dfrac{x^2}{3}+\dfrac{y^2}{2}=1\)}
    {\(\dfrac{x^2}{2}+y^2=1\)}
\end{problem}

\begin{problem}
已知 \(9^m=10\)，\(a=10^m-11\)，\(b=8^m-9\)，则（\quad）

\choices
    {\(a>0>b\)}
    {\(a>b>0\)}
    {\(b>a>0\)}
    {\(b>0>a\)}
\end{problem}

\section{填空题}

\begin{problem}
已知向量 \(\bs a=\left(m,3\right)\)，\(\bs b=\left(1,m+1\right)\). 若 \(\bs a\perp\bs b\)，则 \(m=\)\fillinblank{}
\end{problem}

\begin{problem}
设圆 \(M\) 的圆心在直线 \(2x+y-1=0\) 上，点 \(\left(3,0\right)\) 和 \(\left(0,1\right)\) 均在圆 \(M\) 上，则圆 \(M\) 的方程为\fillinblank{}.
\end{problem}

\begin{problem}
记双曲线 \(C:\dfrac{x^2}{a^2}-\dfrac{y^2}{b^2}=1\)（\(a>0\)，\(b>0\)）的离心率为 \(e\)，写出满足条件“直线 \(y=2x\) 与 \(C\) 无公共点”的 \(e\) 的一个值\fillinblank{}.
\end{problem}

\begin{problem}
已知 \(\triangle ABC\) 中，点 \(D\) 在边 \(BC\) 上，\(\angle ADB=120^\circ\)，\(AD=2\)，且 \(CD=2BD\). 当 \(\dfrac{AC}{AB}\) 取得最小值时，\(BD=\)\fillinblank{}.
\end{problem}

\section{解答题}

\begin{problem}
甲，乙两城之间的长途客车均由 \(A\) 和 \(B\) 两家公司运营，为了解这两家公司长途客车的运行情况，随机调查了甲，乙两城之间的 \(500\) 个班次，得到下面列联表：

\begin{center}
\begin{tabular}{c c c}
\hline
& 准点班次数 & 未准点班次数 \\
\hline
\(A\) & 240 & 20 \\
\(B\) & 210 & 30 \\
\hline
\end{tabular}
\end{center}

\begin{enumerate}
\item 根据上表，分别估计这两家公司甲，乙两城之间的长途客车准点的概率；
\item 能否有 \(90\%\) 的把握认为甲，乙两城之间的长途客车是否准点与客车所属公司有关？
\end{enumerate}

附：\(K^2=\dfrac{n\left(ad-bc\right)^2}{\left(a+b\right)\left(c+d\right)\left(a+c\right)\left(b+d\right)}\)

\begin{center}
\begin{tabular}{c c c c}
\hline
\(P\left(K^2\ge k\right)\) & 0.100 & 0.050 & 0.010 \\
\hline
\(k\) & 2.706 & 3.841 & 6.635 \\
\hline
\end{tabular}
\end{center}
\end{problem}

\begin{problem}
记 \(S_n\) 为数列 \(\{a_n\}\) 的前 \(n\) 项和. 已知 \(\dfrac{2S_n}{n}+n=2a_n+1\).

\begin{enumerate}
\item 证明：\(\{a_n\}\) 是等差数列；
\item 若 \(a_4\)，\(a_7\)，\(a_9\) 成等比数列，求 \(S_n\) 的最小值.
\end{enumerate}
\end{problem}

\begin{problem}
小明同学参加综合实践活动，设计了一个封闭的包装盒，包装盒如图所示：底面 \(ABCD\) 是边长为 \(8\mathrm{~cm}\) 的正方形，\(\triangle EAB\)，\(\triangle FBC\)，\(\triangle GCD\)，\(\triangle HDA\) 均为正三角形，且它们所在的平面都与平面 \(ABCD\) 垂直.

\begin{enumerate}
\item 证明：\(EF\parallel\) 平面 \(ABCD\)；
\item 求该包装盒的容积（不计包装盒材料的厚度）.
\end{enumerate}

\end{problem}

\begin{problem}
已知函数 \(f(x)=x^3-x\)，\(g(x)=x^2+a\)，曲线 \(y=f(x)\) 在点 \(\left(x_1,f(x_1)\right)\) 处的切线也是曲线 \(y=g(x)\) 的切线.

\begin{enumerate}
\item 若 \(x_1=-1\)，求 \(a\)；
\item 求 \(a\) 的取值范围.
\end{enumerate}
\end{problem}

\begin{problem}
设抛物线 \(C:y^2=2px\)（\(p>0\)）的焦点为 \(F\)，点 \(D\left(p,0\right)\)，过 \(F\) 的直线交 \(C\) 于 \(M\)，\(N\) 两点. 当直线 \(MD\) 垂直于 \(x\) 轴时，\(\lvert MF\rvert=3\).

\begin{enumerate}
\item 求 \(C\) 的方程；
\item 设直线 \(MD\)，\(ND\) 与 \(C\) 的另一个交点分别为 \(A\)，\(B\)，记直线 \(MN\)，\(AB\) 的倾斜角分别为 \(\alpha\)，\(\beta\). 当 \(\alpha-\beta\) 取得最大值时，求直线 \(AB\) 的方程.
\end{enumerate}
\end{problem}

\section{选考题：解答题}

\begin{problem}
在直角坐标系 \(xOy\) 中，曲线 \(C_1\) 的参数方程为 \(x=\dfrac{2+t}{6}\)，\(y=\sqrt{t}\)（\(t\) 为参数）；曲线 \(C_2\) 的参数方程为 \(x=-\dfrac{2+s}{6}\)，\(y=-\sqrt{s}\)（\(s\) 为参数）.

\begin{enumerate}
\item 写出 \(C_1\) 的普通方程；
\item 以坐标原点为极点，\(x\) 轴正半轴为极轴建立极坐标系，曲线 \(C_3\) 的极坐标方程为 \(2\cos\theta-\sin\theta=0\)，求 \(C_3\) 与 \(C_1\) 交点的直角坐标，及 \(C_3\) 与 \(C_2\) 交点的直角坐标.
\end{enumerate}
\end{problem}

\begin{problem}
已知 \(a\)，\(b\)，\(c\) 均为正数，且 \(a^2+b^2+4c^2=3\)，证明：

\begin{enumerate}
\item \(a+b+2c\le3\)；
\item 若 \(b=2c\)，则 \(\dfrac{1}{a}+\dfrac{1}{c}\ge3\).
\end{enumerate}
\end{problem}

